% chapters/07/section-01.tex
\section{Kelas \istilah{Sprite} dan Wadah \istilah{Group}}
\label{sec:kelas-sprite}

Pada bab-bab sebelumnya, setiap benda disimpan pada beberapa variabel
terpisah. Cara itu memadai untuk satu atau dua benda, tetapi tidak lagi ketika
jumlahnya bertambah. Bagian ini memperkenalkan cara menyusun benda menjadi satu
kesatuan.

\subsection{Mengapa Benda Perlu Menjadi Satu Kesatuan}

Perhatikan keadaan yang dihadapi pada akhir Bab~VI. Pemain disimpan pada empat
variabel: posisi, kecepatan, \code{rect}, dan penanda menyentuh lantai. Bola
disimpan pada tiga variabel yang serupa. Menambahkan benda ketiga berarti
menambah tiga atau empat variabel lagi, beserta baris pembaruannya.

Keadaan itu membawa tiga kesulitan. Pertama, jumlah variabel bertambah secepat
jumlah benda. Kedua, mudah keliru memasangkan variabel benda yang satu dengan
benda yang lain, karena semuanya bertipe sama. Ketiga, dan yang paling
merepotkan, perintah yang berlaku untuk seluruh benda---misalnya memperbarui
gerak atau menggambar---harus ditulis berulang untuk setiap benda.

Pemrograman berorientasi objek menyelesaikan ketiga kesulitan itu dengan
mengumpulkan data dan perilaku sebuah benda menjadi satu kesatuan yang disebut
objek, dan menyediakan resep pembuatnya yang disebut kelas.

\begin{catatabox}
Apabila Anda belum pernah memakai kelas, tiga hal berikut cukup untuk memulai.
\textbf{Pertama}, kelas adalah resep, sedangkan objek adalah hasilnya: satu
kelas dapat menghasilkan banyak objek. \textbf{Kedua}, fungsi
\code{__init__} dijalankan sekali ketika objek dibuat, dan di sanalah nilai
awalnya disiapkan. \textbf{Ketiga}, kata \code{self} menunjuk objek yang sedang
dikerjakan, sehingga \code{self.rect} berarti ``rect milik objek ini''.
Penjelasan lebih jauh dapat Anda lihat pada buku rujukan \mbox{[R1]},
\mbox{[R2]}.
\end{catatabox}

\subsection{Pola Kelas \istilah{Sprite} pada Pygame}

Pygame menyediakan kelas dasar \code{pygame.sprite.Sprite} yang sudah memuat
keperluan dasar sebuah benda permainan. Tugas Anda adalah menurunkannya dan
mengisi dua hal: gambar benda dan persegi pembatasnya.

\begin{kodekeluaran}
class Item(pygame.sprite.Sprite):
    def __init__(self, x, y):
        super().__init__()
        self.image = pygame.Surface((18, 18))
        self.image.fill(ITEM_COLOR)
        self.rect = self.image.get_rect(center=(x, y))
\end{kodekeluaran}

Tiga baris di dalam \code{__init__} itu perlu dipahami satu per satu.

\begin{itemize}
  \item \code{super().__init__()} menjalankan penyiapan dari kelas dasar.
        Baris itu wajib ada pada setiap kelas turunan \code{Sprite}; tanpa
        itu, benda Anda tidak akan dikenali oleh perkakas Pygame.
  \item \code{self.image} adalah permukaan tempat benda digambar. Namanya harus
        tepat \code{image}, karena perkakas Pygame mencari nama itu.
  \item \code{self.rect} adalah persegi pembatas benda. Namanya pun harus tepat
        \code{rect}, dan nilainya harus berupa \code{Rect}, sebab seluruh
        pemeriksaan tabrakan memakainya.
\end{itemize}

\begin{peringatanbox}
Nama \code{image} dan \code{rect} bukan pilihan Anda, melainkan kesepakatan
Pygame. Apabila keduanya diberi nama lain, program tetap berjalan tetapi
perkakas penggambaran dan pemeriksaan tabrakan tidak menemukan benda Anda, dan
gejalanya membingungkan: benda tidak tergambar, atau tidak pernah dianggap
bertabrakan. Inilah salah satu dari sedikit tempat pada Python yang menuntut
nama tertentu.
\end{peringatanbox}

\subsection{Wadah \istilah{Group}}

Satu kelas menghasilkan banyak benda. Untuk mengurus sekumpulan benda itu,
Pygame menyediakan wadah \code{pygame.sprite.Group}. Wadah menyimpan benda dan
mengerjakan perintah untuk seluruh anggotanya sekaligus. Hubungan keduanya
diperlihatkan pada Gambar~\ref{fig:sprite-group}.

\begin{figure}[htbp]
\centering
% chapters/07/figures/sprite-group.tikz
% Hubungan kelas Sprite dan wadah Group: satu kelas menghasilkan banyak benda,
% dan wadah mengerjakan perintah untuk seluruh bendanya sekaligus.
\begin{tikzpicture}[
  kotak/.style={
    rectangle, rounded corners=2pt, draw=ubPrimary, thick, fill=ubAlt,
    align=center, inner sep=4pt, font=\footnotesize
  },
  wadah/.style={
    rectangle, rounded corners=2pt, draw=ubAccent, thick, fill=white,
    align=center, inner sep=6pt, font=\footnotesize
  },
  benda/.style={circle, draw=ubAccent, thick, fill=ubAccent!25, inner sep=3pt},
  panah/.style={-{Stealth[length=2.4mm]}, thick, draw=ubAccent},
  label/.style={font=\scriptsize, align=center}
]

% Kelas sebagai resep, di sebelah kiri.
\node[kotak, text width=34mm] (kelas) at (0,0)
  {\textbf{kelas \code{Item}}\\[2pt]\scriptsize resep benda: gambar,\\rect, dan perilakunya};

% Tiga benda hasil kelas tersebut, di dalam wadah Group di sebelah kanan.
\node[wadah, text width=48mm, minimum height=20mm] (grup) at (7.6,0) {};
\node[anchor=south, font=\scriptsize\bfseries, ubAccent] at (7.6,1.15)
  {wadah \code{Group}};

\node[benda] at (6.4,-0.1) {};
\node[benda] at (7.6,-0.1) {};
\node[benda] at (8.8,-0.1) {};

\draw[panah] (kelas.east) -- node[label, above, yshift=1pt]
  {membuat tiga benda} (grup.west);

% Keterangan di bawah diagram, pada barisnya sendiri.
\node[label, anchor=north, text width=124mm] at (3.8,-1.7)
  {\code{update()} dan \code{draw()} diberikan kepada wadah, lalu wadah
   mengerjakannya untuk setiap benda di dalamnya. Perintah tidak perlu ditulis
   satu per satu.};

\node[label, anchor=north, text width=124mm] at (3.8,-2.7)
  {Satu kelas dipakai berulang kali. Menambah benda baru cukup dengan membuat
   satu benda lagi dan memasukkannya ke dalam wadah.};
\end{tikzpicture}

\caption{Kelas menghasilkan benda, wadah mengurus sekumpulan benda}
\label{fig:sprite-group}
\end{figure}

Tabel~\ref{tab:sprite-group} memuat anggota kedua kelas yang paling sering
dipakai. Seluruhnya disediakan Pygame, sehingga tidak perlu ditulis sendiri
\cite{pygame-sprite-mask}.

\begin{table}[htbp]
\centering
\caption{Anggota \code{Sprite} dan \code{Group} yang paling sering dipakai}
\label{tab:sprite-group}
{\setlength{\tabcolsep}{4pt}\footnotesize%
\begin{tabularx}{\textwidth}{@{}L{4.6cm} Y@{}}
\toprule
\headrow \textbf{Penulisan} & \textbf{Arti} \\
\midrule
\code{sprite.image}, \code{sprite.rect} &
Gambar dan persegi pembatas benda. Kedua nama itu wajib dipakai. \\
\addlinespace
\code{sprite.update(dt, ...)} &
Perilaku benda setiap bingkai. Isinya Anda tentukan sendiri; argumen tambahan
diteruskan oleh wadah. \\
\addlinespace
\code{sprite.kill()} &
Mengeluarkan benda dari seluruh wadah yang memuatnya. Berguna ketika benda
diambil atau dihancurkan. \\
\addlinespace
\code{group.add(benda)} &
Memasukkan benda ke dalam wadah. Satu benda boleh berada pada beberapa wadah
sekaligus. \\
\addlinespace
\code{group.update(dt, ...)} &
Memanggil \code{update} pada seluruh anggotanya, dengan argumen yang sama. \\
\addlinespace
\code{group.draw(surface)} &
Menggambar seluruh anggotanya ke permukaan yang diberikan. \\
\addlinespace
\code{for benda in group:} &
Menelusuri seluruh anggota wadah, sebagaimana menelusuri daftar biasa. \\
\bottomrule
\end{tabularx}}
\end{table}

\begin{tipsbox}
Satu benda boleh berada pada beberapa wadah sekaligus. Cara
itu dipakai pada program bab ini: seluruh bongkahan dikumpulkan di wadah
\code{platforms}, dan apabila nanti perlu menggambar seluruh isi permainan
dengan satu perintah, bongkahan yang sama dapat dimasukkan pula ke wadah
\code{all_sprites}. Wadah hanyalah kumpulan rujukan, bukan salinan.
\end{tipsbox}
